\documentclass[12pt, aspectratio=169]{beamer}

\usepackage[utf8]{inputenc}
\usepackage[T2A]{fontenc}
\usepackage[russian]{babel}

\usepackage{graphicx}
\usepackage{booktabs}
\usepackage{array}
\usepackage{tikz}
\usetikzlibrary{arrows.meta, positioning}

\usetheme{Madrid}

\definecolor{ogublue}{RGB}{30, 60, 140}
\definecolor{ogugold}{RGB}{200, 160, 30}
\definecolor{ogulightblue}{RGB}{220, 230, 250}

\setbeamercolor{palette primary}{bg=ogublue, fg=white}
\setbeamercolor{palette secondary}{bg=ogublue!80, fg=white}
\setbeamercolor{palette tertiary}{bg=ogublue!60, fg=white}
\setbeamercolor{frametitle}{fg=white, bg=ogublue}
\setbeamercolor{block title}{fg=white, bg=ogublue}
\setbeamercolor{block body}{fg=black, bg=ogulightblue}
\setbeamercolor{block title alerted}{fg=white, bg=red!70!black}
\setbeamercolor{block body alerted}{fg=black, bg=red!10}
\setbeamercolor{itemize item}{fg=ogublue}
\setbeamercolor{itemize subitem}{fg=ogugold}

\setbeamertemplate{navigation symbols}{}
\setbeamertemplate{itemize item}{\textbullet}
\setbeamertemplate{itemize subitem}{--}

\title[ФМФ ОГУ]{\textbf{История Физико-математического факультета ОГУ}}
\subtitle{Орловский государственный университет им.\,И.\,С.\,Тургенева}
\author[Яшин К.\,А.]{Яшин К.\,А.}
\date{22 майя 2026}

\begin{document}

{\setbeamertemplate{footline}{}
\begin{frame}
  \titlepage
\end{frame}}

\begin{frame}{Введение}
  \begin{columns}[T]
    \column{0.55\textwidth}
      \begin{block}{}
        Физико-математический факультет ОГУ~--- \textbf{один из старейших}
        в университете, стоявший у истоков его создания.
      \end{block}
      \vspace{0.4em}
      \begin{block}{}
        Факультет имеет богатую историю
        подготовки учителей физики,
        математики и информатики, сочетая
        фундаментальное образование с
        актуальными научными
        исследованиями в области физики,
        математики и IT-технологий.
      \end{block}
    \column{0.4\textwidth}
  \begin{center}
    \includegraphics[width=0.9\textwidth]{logotype.jpg}\\[0.3em]
    {\small\textcolor{ogublue}{\textit{Эмблема факультета}}}
  \end{center}
  \end{columns}
\end{frame}

\begin{frame}{История основания}
  \begin{center}
  \begin{tikzpicture}[
    box/.style={draw=ogublue, fill=ogulightblue, rounded corners=5pt,
                minimum width=2.6cm, minimum height=1.1cm,
                align=center, line width=1.2pt, font=\small},
    arrow/.style={->, >=Stealth, ogublue, line width=1.5pt}
  ]
    \node[box] (n1) at (-3.6, 0) {\textbf{5.08.1931}\\\footnotesize Основание\\института};
    \node[box] (n2) at (0,    0) {\textbf{16.10.1931}\\\footnotesize Физ.-технич.\\факультет};
    \node[box] (n3) at (3.6,  0) {\textbf{1932}\\\footnotesize Физ.-мат.\\факультет};
    \draw[arrow] (n1)--(n2);
    \draw[arrow] (n2)--(n3);

    \node[draw=ogugold, fill=ogugold!12, rounded corners=4pt,
          minimum width=10.0cm, align=center, below=0.65cm of n2,
          font=\itshape\small, line width=1pt]
      {<<Один из старейших факультетов ОГУ>>};
  \end{tikzpicture}
  \end{center}
\end{frame}

\begin{frame}{Развитие факультета}
  \begin{columns}[T]
    \column{0.47\textwidth}
      \begin{alertblock}{1950-е годы}
        Функционирование вечерних отделений; активный рост
        числа студентов и преподавателей.
      \end{alertblock}
      \vspace{0.5em}
      \begin{block}{Начало 1980-х --- расширение}
        \begin{itemize}
          \item Кафедра общей физики
          \item Кафедра теоретической физики
          \item Рост научных публикаций
        \end{itemize}
      \end{block}
    \column{0.48\textwidth}
  \begin{block}{Владимир Федотович Пивень}
    \begin{columns}[T]
      \column{0.38\textwidth}
        \includegraphics[width=\textwidth, height=2.8cm, keepaspectratio]{Владимир.Ф.П.png}
      \column{0.58\textwidth}
        д.\,ф.-м.\,н., профессор,\\
        зав. кафедрой общей физики.\\[0.4em]
        \begin{itemize}
          \item Проработал \textbf{более 46~лет}
          \item Основатель научной школы
          \item Подготовил десятки учёных
        \end{itemize}
    \end{columns}
  \end{block}
  \end{columns}
\end{frame}

\begin{frame}{Современность}
  \begin{columns}[T]
    \column{0.48\textwidth}
      \begin{block}{Структура (2025)}
        \textbf{5 кафедр факультета:}
        \begin{itemize}
          \item Алгебры и геометрии
          \item Информатики
          \item Математического анализа
          \item Экспер. и теор. физики
          \item Математики и ПИТ
        \end{itemize}
        \medskip
        \textbf{94} бюдж. места на 2025 год \newline от \textbf{128\,240 руб/год}
      \end{block}
    \column{0.48\textwidth}
      \begin{block}{Востребованные профессии}
        \begin{itemize}
          \item Системный аналитик
          \item Руководитель IT-проектов
          \item Разработчик приложений
          \item Инженер прикладной сферы
        \end{itemize}
      \end{block}
      \vspace{0.3em}
      \begin{alertblock}{Научная конференция}
        \footnotesize Всерос. конф. <<Современные проблемы
        физико-математических наук>>\\
        (ML, анализ данных, методика преподавания)
      \end{alertblock}
  \end{columns}
\end{frame}

\begin{frame}{Кафедры факультета}
  \begin{center}
  \renewcommand{\arraystretch}{1.38}
  \begin{tabular}{>{\bfseries\color{ogublue}}l l l}
    \toprule
    \normalfont\bfseries Кафедра & \bfseries Заведующий & \bfseries Контакты \\
    \midrule
    Мат. анализа и методики & Чаплыгина Е.\,В. & matdiff@yandex.ru \\[2pt]
    Алгебры и мат. методов  & Строев С.\,П.    & 8\,(4862)\,75-24-25 \\[2pt]
    Экспер. и теор. физики  & Марков О.\,И.    & kafedrafiziki.ogu@mail.ru \\[2pt]
    Информатики             & Дорофеева В.\,И. & kafedra.informatiki@ogu.ru \\[2pt]
    Математики и ПИТ        & Авдеев И.\,Ф.    & 8\,(4862)\,43-17-40 \\
    \bottomrule
  \end{tabular}
  \end{center}
  \vspace{0.4em}
  \begin{center}
    \textcolor{ogublue}{\textbf{Адрес:}}\;
    г.\,Орёл, ул.\,Комсомольская, д.\,95 (корпус \textnumero 1)
  \end{center}
\end{frame}

\begin{frame}{Заключение}
  \begin{block}{Итоги}
    Физико-математический факультет ОГУ за \textbf{более 90~лет} своей истории
    прошёл путь от небольшого отделения до крупного современного факультета
    с пятью кафедрами и богатой научной школой.
  \end{block}
  \vspace{0.5em}
  \begin{columns}[T]
    \column{0.32\textwidth}
      \begin{alertblock}{\centering Основан}
        \centering\Large\textbf{1931}
      \end{alertblock}
    \column{0.32\textwidth}
      \begin{alertblock}{\centering Кафедр}
        \centering\Large\textbf{5}
      \end{alertblock}
    \column{0.32\textwidth}
      \begin{alertblock}{\centering Бюдж. мест}
        \centering\Large\textbf{94}
      \end{alertblock}
  \end{columns}
  \vspace{0.7em}
  \begin{center}
    \textcolor{ogublue}{\large\textbf{Спасибо за внимание!}}
  \end{center}
\end{frame}

\end{document}